%============================================================ 7. 고찰 및 결론
\section{고찰 및 결론}
\label{sec:discussion}

\subsection{결과가 말하는 것}

\S\ref{sec:results}의 결과를 실무 판단으로 옮기면 세 가지다. 기동 계층은 조건 없이 넣는다 ---
언어모델 없이도 같은 방어 성과를 더 적은 자원으로 내며 외부 의존이 없다. 지휘관 계층은
\emph{능력 임계}를 넘는 모델에만 켠다. 본 실험에서 그 경계는 로컬 14B와 클라우드 모델 사이였고
로컬 7B는 코드 휴리스틱보다 유의하게 나빴다(\S\ref{sec:res_threshold}). 두 계층의 기여는
가산적이므로(\S\ref{sec:res_g3}) 지휘관 모델 교체와 기동 정책 재학습을 독립적으로 진행할 수
있다.

지휘관의 성능을 설명하는 것은 커버리지가 아니라 재계획 간 계획 안정성이라는
관찰(\S\ref{sec:res_mechanism})은 본 문제에 국한되지 않는다. 주기적으로 재계획하는 계층
구조에서 상위 결정의 변경은 하위 계층에 \emph{전환 비용}을 발생시킨다 --- 그물벽을 절반쯤
전개한 방어정의 담당이 바뀌면 그때까지의 이동과 그물이 회수되지 않는다. 상위 계층이 매
재계획마다 순간 최적을 추구하면 이 비용이 누적되어 국소적으로 나은 배정이 전역적으로 손해가
된다. 언어모델은 매 호출이 독립적이라 이 일관성이 구조적으로 보장되지 않으므로, 직전 배정을
상태에 포함시키고 담당 무리를 방위로 식별하게 한 것이 그 대응이다. 따라서 언어모델을 계획
계층에 둘 때는 계획 품질을 국소 최적성이 아니라 \emph{시간에 걸친 일관성}으로 평가해야 하며,
프롬프트와 상태 표현이 그 일관성을 지원해야 한다.

가장 큰 개선이 총 선회량($\dz = -3.23$, \S\ref{sec:res_effectsize})에서 나타난 것은 점수맵
표현의 성질과 직결된다. 경유점이 격자에 고정되므로 연속 회귀의 미세 진동이 구조적으로 배제되어
궤적이 부드러워지며, 선회 감소는 실제 선박에서 연료·시간·기동 여유로 환산된다. 포획당 그물
소모의 감소($-45\,\%$) 역시 유효 마스크가 배정 회랑과 사거리 제약을 행동 분포에 직접 부과해
무의미한 위치의 그물벽이 애초에 표현되지 않기 때문이다.

배치 관점에서 로컬 14B는 클라우드 대비 10.6배 느리면서 성능도 낮았고(0.811 대 0.904) 호출의
4.2\,\%가 결정 주기를 초과하였으므로(\S\ref{sec:res_latency}, Table~\ref{tab:latency_tail}),
비동기 설계는 로컬 배치의 전제 조건이다.

\subsection{한계와 향후 과제}
\label{sec:limitations}

\textbf{시뮬레이션 한정.} 모든 결과는 시뮬레이션에서 얻었으며, 그물의 유체 거동(선분 모형),
추진기 얽힘 판정(접촉 즉시 무력화), 관측 완전성(무오차 위치)이 단순화되어 있다. 이 가운데
관측 불확실성은 배정이 의존하는 무리 구성을 흔들어 churn 을 직접 키우므로, 점수맵 정책보다
지휘관 계층에 더 크게 작용할 것이다. 지연 측정이 GPU 공유 구성의 교란을 받는 점은
Table~\ref{tab:latency_tail}의 주석에 밝혔다.

\textbf{폐쇄형 모델 의존.} 가장 좋은 지휘관은 상용 API 의 폐쇄형 모델이라 그 수치의 장기
재현은 보장되지 않고, 능력 임계의 위치(\S\ref{sec:res_threshold})도 시점 의존적이다. 근거
원문을 저장하지 않아 판단 근거의 사후 검토가 제한된다. 기동 계층 단독 조건은 언어모델을
쓰지 않으므로 이 위험과 무관하게 재현된다.

\textbf{규모와 일반화.} 실험은 적 10척 대 방어정 3척의 고정 규모이며, 방어정 수가 늘어
배정 조합이 급증할 때 능력 임계가 어디로 이동하는지는 측정되지 않았다. 세 공격 포메이션은
규칙 기반이고, 그물벽 배치 패턴을 학습해 회피하는 적응적 적은 범위 밖이다.

\subsection{결론}
\label{sec:conclusion}

시드 대응 $2\times2$ 설계로 두 계층의 기여를 분리 측정한 720
에피소드에서 세 가지가 얻어졌다. 첫째, 두 계층의 기여는 \textbf{가산적}이다 --- 10개 지표 중
9개에서 상호작용의 신뢰구간이 0을 포함하였고, 본 논문은 이를 상승효과로 포장하지 않았다.
둘째, \textbf{포획률은 측정한 지표 중 가장 둔하다}. 효과크기로 총 선회량($-3.23$)과 포획당
그물 소모($-2.09$)의 개선이 포획률($+0.44$)보다 훨씬 크고, 다중비교 보정을 걸면 10개 지표 중
8개가 견디는데 포획률과 돌파 수만 유의성을 잃는다. 제안 체계의 실질적 이득은 ``더 많이
막는다''보다 ``같은 정도로 막으면서 훨씬 적은 자원과 위험으로 해낸다''에 있다. 셋째,
\textbf{언어모델 지휘관에는 능력 임계가 존재하며}, 그 성능을 설명하는 것은 배정 커버리지가
아니라 재계획 간 계획 안정성이다. 가장 낮은 성능의 모델이 오히려 커버리지는 가장 높았다는
관측은 주기적 재계획 구조에서 최적성보다 일관성이 중요하다는 설계 원리를 시사한다.

검증은 시뮬레이션에 한정되며, 실해역 이식과 관측 불확실성 하에서의 계획 안정성은 향후 과제로
남는다.
